\section{Entire original curve FULL values across every genuine history prefix}

\begin{theorem}[The actual curve-reference formula rebases across every genuine prefix]
\label{hir:original-modeled-curve-full-arbitrary-prefix-rebase}
\uses{hir:original-modeled-untested-curve-full-value,bcalc:inv-domain-stable}
\lean{Hironaka.HironakaOriginalCurveFullRebase.modeled_curve_full_append_rebase}
\leanok
Let $m,K$ be original marked ideals on the SAME ambient manifold. Suppose
that every genuine TEST $R$ of $m$ has actual transformed reference $K_R$
in DOMAIN$(1)$. For every actual prefix $H$ and tail $S$ whose concatenation
is a genuine original TEST of $m$, every point in the current reference
cosupport, and every $k\ge0$, the unchanged original FULL value at depth
$k+1$ satisfies
\[
 \operatorname{FULL}_{k+1}(K,H\mathbin{\cdot}S;y)
 =\operatorname{FULL}_{k+1}(K_H,S;y).
\]
Neither $H$ nor $S$ is required to be empty or canonical. No TEST of the
reference $K$, or numerical rebasing equality, is a hypothesis.
\end{theorem}
\begin{proof}
The true TEST concatenation equivalence supplies the actual prefix TEST
of $m$ and tail TEST of $m_H$. The genuine reference-domain hypothesis
therefore gives DOMAIN$(1)$ for $K_H$. For ANY actual tail TEST $R$ of
$m_H$, its real concatenation $H\mathbin{\cdot}R$ is a TEST of $m$.
Applying the hypothesis and the exact whole controlled-transform append
identity proves that $(K_H)_R$ is in DOMAIN$(1)$. This constructs the
entire reference-domain family for the actual tail, rather than assuming
a tail TEST of $K_H$.

Apply the proved explicit actual curve FULL formula first to $K$ on
$H\mathbin{\cdot}S$ and then to $K_H$ on $S$. Each formula has two literal
branches. At order zero it is the original monomial invariant of the
ENTIRE current record; at positive order its second coordinate is the
actual zero-level special value with top lower remainder. The whole
controlled-transform append identity identifies these two current marked
records, hence their nonmonomial order, monomial invariant, branch test
and positive branch. The two displayed FULL expressions are consequently
equal at EVERY positive depth. No historical canonicality is used.
\end{proof}

\begin{theorem}[Only original DOMAIN one internally supplies arbitrary-prefix numerical rebasing]
\label{hir:original-curve-full-arbitrary-prefix-rebase}
\uses{hir:original-modeled-curve-full-arbitrary-prefix-rebase,bcalc:inv-domain-stable}
\lean{Hironaka.HironakaOriginalCurveFullRebase.original_curve_full_append_rebase}
\leanok
For ONLY original DOMAIN$(1)$ data and any genuine original TEST
$H\mathbin{\cdot}S$, the preceding entire original FULL rebasing holds
at every current cosupport point and EVERY positive depth. No
reference-domain family is supplied by the caller.
\end{theorem}
\begin{proof}
Choose the original reference to be the original marked ideal itself.
Actual preservation of DOMAIN$(1)$ along EVERY genuine original TEST
internally supplies the whole reference-domain family required by the
first theorem. Apply it to the given actual concatenated history and
current cosupport point. The reference and its transformed tail are
precisely the original marked ideal and its true prefix transform, so
this is the claimed original-only equality.
\end{proof}

\begin{theorem}[Actual positive member references rebase without a member TEST]
\label{hir:original-positive-member-curve-full-arbitrary-prefix-rebase}
\uses{hir:original-modeled-curve-full-arbitrary-prefix-rebase,hir:member-tail-actual-geometry}
\lean{Hironaka.HironakaOriginalCurveFullRebase.original_positive_member_curve_full_append_rebase}
\leanok
For original DOMAIN$(2)$ data, a nonempty actual set of original boundary
positions selects its true member-remainder reference $K$. For EVERY
genuine TEST $H\mathbin{\cdot}S$ of the ORIGINAL marked ideal and every
point in the current $K$ cosupport, the SAME arbitrary-prefix original
FULL rebasing holds at every depth $k+1$. No TEST of $K$ or of its prefix
transform is assumed.
\end{theorem}
\begin{proof}
The original marked ideal is well formed, has positive mark and carrier
dimension at most two. The independently proved actual member-remainder
geometry therefore puts the selected reference in DOMAIN$(1)$ after
EVERY genuine TEST of that original ideal. Nonemptiness of the selected
actual positions is exactly the dimension-drop premise of that theorem.
This internally supplies the full geometric reference-domain family of
the first theorem. Apply that theorem to the SAME original history,
actual selected reference and current cosupport point. It proves the
unchanged original numerical equality, while retaining every real centre
and every empty position. No member TEST, canonicality or desired value
is inferred or supplied.
\end{proof}

This is a numerical rebasing producer for actual curve references. It does
not prove higher-carrier arbitrary-history FULL rebasing or a global
resolution scheduler.
